\section[Stage 039]{Stage 039: First Non-Flat \texorpdfstring{$U$}{U} Doublet and Direction Splitting}
\label{stage:039}

\stagefield{Part / anchor}{Part III; primary anchor \resultanchor{MTDC-T6}.}
\claimstatus{\StatusExactClosure{} inside the first split--\(U\) continuum extension.}
\stagefield{Purpose}{Stage~039 tests the simplifying assumption that the source, support, and loading directions are collinear.  The first non-flat \(U\) doublet preserves scalar placement but generically breaks direction collinearity.}
\stagefield{Inputs}{The Stage~038 flat--\(U\) placement map and the first axial excitation of the internal \(U\) sector.}

\stagefield{Verification}{SymPy audit: \StageFile{scripts/moving_throat_pde_stage039_split_u_sector_sympy_audit.py}; transcript: \StageFile{scripts/output/moving_throat_pde_stage039_split_u_sector_sympy_audit.txt}.  Mathematica audit: \StageFile{mathematica/moving_throat_pde_stage039_split_u_sector_mathematica_audit.wl}; transcript: \StageFile{mathematica/output/moving_throat_pde_stage039_split_u_sector_mathematica_audit.txt}.}

\subsection*{Derivation}

Introduce
\begin{equation}
\label{eq:app-stage039-U-doublet}
K_{U0}=K_U,
\qquad
K_{U1}=K_U(1+\delta_U),
\qquad
\delta_U=\frac{\pi^2T_U}{L^2K_U}.
\end{equation}
The scalar placement data become
\begin{equation}
\label{eq:app-stage039-split-placement}
\boxed{
\delta_{\rm split}
=\frac{\delta_0+\epsilon_\eta\delta_U/(1+\delta_U)}{1-\epsilon_\eta},
\qquad
\epsilon_{W,{\rm split}}
=\epsilon_W\left[1-\frac{2}{11}\frac{\delta_U}{1+\delta_U}\right]}.
\end{equation}
The directional deformation is carried by
\begin{equation}
\label{eq:app-stage039-RU}
\boxed{
R_U=\frac{1+\rho_0/(1+\delta_U)}{1+\rho_0}}.
\end{equation}
The exact direction-splitting invariant is
\begin{equation}
\label{eq:app-stage039-Ddir}
\boxed{
D_{\rm dir}=-\frac{\kappa_0\kappa_1g_W\rho_0\delta_U}{1+\delta_U}}.
\end{equation}
Therefore
\begin{equation}
\label{eq:app-stage039-collinearity}
\boxed{
D_{\rm dir}=0
\iff
\delta_U=0\ \text{or}\ \rho_0=0}.
\end{equation}
On the constructive branch, \(\rho_0>0\) and \(\delta_U>0\) imply a genuine direction split.

\stagefield{Output}{The split placement law \eqref{eq:app-stage039-split-placement}, the direction factor \eqref{eq:app-stage039-RU}, and the collinearity theorem \eqref{eq:app-stage039-collinearity}.}
\stagefield{Downstream use}{Stage~040 packages this direction split into generalized selected-branch functions; Stages~044--046 test whether the resulting branch can still satisfy the target.}

